% !TeX root = ../thesis.tex
% Kapitel 1
\chapter{Einleitung}
\label{cha:einleitung}

\section{Motivation und Kontext}
\label{sec:motivation}
Im vergangenen Jahr habe ich bei der \hsmafirma{} mit einem Team gearbeitet, das ein Webanwendung entwickelt, welche unteranderem das Risikomanagement ermöglichen und erleichtern soll.
Mit dieser Anwendung können Sicherheitsrisiken erfasst, bewertet und überwacht werden.
Dabei habe ich mich häufig mit der Frage beschäftigt, wie wir Risikobewerter so unterstützen können, dass ihre Bewertungen konsistenter und effizienter werden.

Besonders Verbesserungspotential sehe ich in der Wiederverwendung von bereits bewerteten Risiken.
Ein relevanter Aspekt der Risikobewertung ist die Zuordnung von Kontrollen zu Risiken.
Sie ist ein zentraler Bestandteil jeder Risikobewertung und wird derzeit für jedes neue Risiko von Grund auf durchgeführt.
Das ist vor allem dann ineffizient, wenn ein ähnliches Risiko bereits bewertet wurde und die zugehörigen Kontrollen bekannt sind.

Innerhalb dieser Arbeit will ich daher untersuchen, ob es möglich ist, zu einem neuen Sicherheitsrisiko allein anhand seiner Attribute die ähnlichsten bereits bewerteten Risiken zu finden, sodass deren Kontrollfamilien als Wiederverwendungsempfehlung vorgeschlagen werden können.

\section{Problemstellung}
\label{sec:problem}
Die Beantwortung dieser Forschungsfrage wird durch mehrere Eigenschaften der Daten und des Anwendungsfalls erschwert.

\paragraph{Gemischte Attribute:}
Ein Risiko wird durch kategoriale, numerische und textuelle Attribute beschrieben.
Während sich die ersten beiden Typen mit etablierten Distanzmaßen vergleichen lassen, müssen die Freitexte zunächst in eine vergleichbare Darstellung überführt werden.

\paragraph{Fehlende Kategorisierung von Kontrollen:}
Die Kontrollen eines Risikos liegen nur als Freitext vor.
Eine Kategorisierung, auf der eine Empfehlung aufbauen könnte, ist nicht vorhanden.
Die Risikobewerter folgen zwar einem internen Schema, dieses ist jedoch nicht in den Daten abgebildet und steht mir als Quelle nicht zur Verfügung.
Die Kontrollen müssen daher zunächst zu Kontrollfamilien zusammengefasst werden, auf deren Ebene dann verglichen und empfohlen wird.

\paragraph{Mehrere Kontrollfamilien je Risiko:}
Die Kontrollen eines Risikos verteilen sich in der Regel auf mehrere Kontrollfamilien.
Anders als bei einer klassischen Klassifikation, bei der jedem Objekt genau eine Klasse zugeordnet wird, muss unser Verfahren eine Menge passender Kontrollfamilien empfehlen.

\paragraph{Ungleiche Verteilung:}
Die verschiedenen Kontrollfamilien sind ungleichmäßig verteilt.
Für seltene Kontrollfamilien liegen kaum Vergleichsfälle vor, was eine belastbare Empfehlung erschwert.

\paragraph{Unvollständige Entwürfe:}
Eine Empfehlung ist vor allem früh im Bewertungsprozess nützlich.
In der Entwurfsphase sind jedoch viele Attribute noch nicht bekannt, sodass die Ähnlichkeitssuche auch mit unvollständigen Daten durchgeführt funktionieren muss.

\paragraph{Vertraulichkeit:}
Die Daten, die ich für diese Arbeit verwende, sind vertraulich.
Externe Dienste sind daher nicht zulässig, sodass alle Verfahren lokal auf den Daten ausgeführt werden müssen.

% Einschätzung statt Confidence?
\section{Forschungsfrage und Hypothese}
\label{sec:forschungsfrage}
\begin{quote}
  \itshape Welche Methoden eignen sich, um zu einem neuen Sicherheitsrisiko
  allein anhand seiner Attribute die ähnlichsten bereits bewerteten Risiken zu
  finden, sodass deren Kontrollfamilien mit einer belastbaren Einschätzung zur Wiederverwendung empfohlen werden können?
\end{quote}

Hypothese: Weisen Sicherheitsrisiken in ihren relevanten Attributen ähnliche
Werte auf, lassen sich die Kontrollfamilien bekannter Risiken auf neue, ähnliche
Risiken übertragen.

Zur Beantwortung der Forschungsfrage werden folgende Teilfragen untersucht:

\begin{enumerate}
  \item\label{tf: attribute} Verbessert eine Selektion innerhalb der vorgefilterten Attribute die Empfehlung von Kontrollfamilien und wie stabil ist die Auswahl über die Durchgänge?
  \item\label{tf: methoden} Übertreffen die untersuchten Ähnlichkeitsverfahren ein einfaches Referenzverfahren bei der Empfehlung von Kontrollfamilien und welches Verfahren erzielt dabei die besten Ergebnisse?
  \item\label{tf: robustheit} Wie robust sind die Empfehlungen bei fehlenden Attributen und bleibt ihre Einschätzung belastbar?
  % \item\label{tf: kausalitaet} Bleibt der Zusammenhang zwischen Attributähnlichkeit und Kontrollfamilien bestehen, wenn Störfaktoren wie Org-Einheit oder Bewerter berücksichtigt werden?
\end{enumerate}

\section{Beitrag und Abgrenzung}
\label{sec:beitrag}
Der Beitrag dieser Arbeit ist ein systematischer Vergleich von Ähnlichkeitsverfahren zur Empfehlung von Kontrollfamilien auf realen Daten eines Security Risk Registers.
Dazu wurde eine Evaluationspipeline entwickelt, die Attributselektionen, Ähnlichkeitsverfahren und Simulationen von Entwurfsphasen in mehreren Durchgängen anwendet und dabei ausschließt, dass Informationen aus der Testmenge in die Trainingsmenge gelangen.
Die Ergebnisse zeigen, welche Verfahren sich eignen, wie robust sie in der Entwurfsphase sind und wie belastbar ihre Einschätzung ist.

Das Verfahren dient ausschließlich als Entscheidungshilfe für Risikobewerter und soll keine Kontrollen vorschreiben.

\section{Aufbau der Arbeit}
\label{sec:aufbau}
Kapitel~\ref{cha:grundlagen} liefert einen Einblick in die Risikobewertung und erläutert die Grundlagen zu Attributselektion, Ähnlichkeitsmaßen, Empfehlungsverfahren und deren Bewertung.
Kapitel~\ref{cha:verwandt} stellt verwandte Arbeiten zu Vergleichsverfahren und Empfehlungssystemen vor.
Dazu wird die vorliegende Arbeit in den Kontext der Forschung eingeordnet.
Kapitel~\ref{cha:methodik} beschreibt die verwendeten Daten, ihre Aufbereitung und die Evaluationspipeline.
In Kapitel~\ref{cha:evaluation} werden die Experimente durchgeführt, die Ergebnisse dargestellt und im Hinblick auf die Teilfragen diskutiert.
Kapitel~\ref{cha:fazit} fasst die Arbeit zusammen, beantwortet die Forschungsfrage und gibt einen Ausblick auf mögliche Weiterentwicklungen.